\subsection{标准正交族}

定义 3. --- 在准希尔伯特空间中，称向量族 \((e_{i})_{i\in\mathbf{I}}\) 为正交的，如果 \(e_{i}\) 与 \(e_{k}\)（对所有 \(i\neq k\)）正交；称其为标准正交的，如果此外 \(\|e_{i}\|=1\) 对所有 \(i\in I\) 成立。

E 的子集 S，若由 S 到其自身的恒等映射所定义的族是标准正交的，则称之为标准正交集。若 \((e_{i})_{i\in I}\) 是标准正交族，则映射 \(i\mapsto e_{i}\) 是单射；于是我们可以不加区分地谈论标准正交族或标准正交集。

若 \((e_i)_{i\in I}\) 是标准正交族，则完备的一维向量子空间 \(\mathbf{D}_i = \mathbf{K}e_i\) 两两正交。对每个 \(x\in E\)，\(x\) 在 \(\mathbf{D}_i\) 上的正交投影是 \(\lambda_i e_i\)，其中 \(\langle e_i|x - \lambda_i e_i\rangle = 0\)，由此得 \(\langle e_i|x\rangle = \lambda_i\langle e_i|e_i\rangle = \lambda_i\)。将 No.~2 的结果应用于诸子空间 \(\mathbf{D}_i\)，便得到下列命题：

命题 3. --- 在分离的准希尔伯特空间 E 中，每个标准正交族都是拓扑无关的。

我们注意到，这一性质可由拓扑无关族的刻画（IV, p.~1 与 II, p.~43, 推论 2）立即得出，因为按 K 等于 R 或 C，E 的对偶分别与 E 的完备化或 E 的共轭空间等同（V, p.~17, 注）。

命题 4. --- 设 E 为分离的准希尔伯特空间，\((e_{i})_{i\in I}\) 为 E 中的标准正交族，V 为 E 中由诸 \(e_{i}\) 生成的闭向量子空间。

\begin{enumerate}
\def\labelenumi{\arabic{enumi})}
\tightlist
\item
  对每个 \(x \in E\)，我们有
\end{enumerate}

\[
\sum_ {i \in \mathbf {I}} \left| \langle e _ {i} | x \rangle \right| ^ {2} \leqslant \| x \| ^ {2}\tag{1}
\]

（Bessel 不等式）；这里，所有 \(i \in \mathbf{I}\) 中使 \(\langle e_i|x\rangle \neq 0\) 者所成的集是可数的。此外，下列条件等价：a) \(x \in \mathbf{V}\)；b) \(\|x\|^2 = \sum_{i \in \mathbf{I}} |\langle e_i|x\rangle|^2\)；c) 族 \(\langle e_i|x\rangle\) . \(e_i\) 在 E 中可和，且 \(x = \sum_{i \in \mathbf{I}} \langle e_i|x\rangle\) . \(e_i\)。

\begin{enumerate}
\def\labelenumi{\arabic{enumi})}
\setcounter{enumi}{1}
\item
  若 \(\mathbf{V}\) 完备，则由所有 \(\langle e_i|x\rangle\) . \(e_i\) 组成的族在 \(\mathbf{E}\) 中（对所有 \(x\in\mathbf{E}\)）可和，且 \(\sum_{i\in\mathbf{I}}\langle e_i|x\rangle\) . \(e_i=p_{\mathbf{V}}(x)\)，\(\sum_{i\in\mathbf{I}}|\langle e_i|x\rangle|^2=\|p_{\mathbf{V}}(x)\|^2\)。
\item
  假设 V 完备。对每个标量族 \((\lambda_{i})_{i\in\mathbf{I}}\)（满足 \(\sum_{i\in I}|\lambda_{i}|^{2}<+\infty\)），存在唯一的点 \(x\in V\)，使得 \(\langle e_{i}|x\rangle=\lambda_{i}\) 对所有 \(i\in I\) 成立。若 \((\mu_{i})_{i\in I}\) 是第二个满足 \(\sum_{i\in I}\left|\mu_{i}\right|^{2}<+\infty\) 的标量族，且 \(y\in V\) 满足 \(\langle e_{i}|y\rangle=\mu_{i}\)（对所有 \(i\in I\)），则 \(\langle x|y\rangle=\sum_{i\in I}\overline{\lambda}_{i}\mu_{i}\)。
\end{enumerate}

命题 5. --- 设 \((e_i)_{i \in I}\) 为分离的准希尔伯特空间 E 中的标准正交族。下列性质等价：

\begin{enumerate}
\def\labelenumi{\alph{enumi})}
\item
  族 \((e_i)\) 是完全的；
\item
  对每个 \(x \in E\)，族 \(\langle e_{i}|x\rangle\) . \(e_{i}\) 在 E 中可和，且 \(x = \sum_{i \in I} \langle e_{i}|x\rangle\) . \(e_{i}\)；
\item
  对每个 \(x \in E\)，
\end{enumerate}

\[
\left\| x \right\| ^ {2} = \sum_ {i \in \mathbf {I}} \left| \langle e _ {i} | x \rangle \right| ^ {2}\tag{2}
\]

（Parseval 等式）。

当 E 是希尔伯特空间时，这些条件也等价于：

\begin{enumerate}
\def\labelenumi{\alph{enumi})}
\setcounter{enumi}{3}
\tightlist
\item
  关系式 \(\langle e_i|x\rangle = 0\)（对所有 \(i\in \mathrm{I}\)）蕴含 \(x = 0\)。
\end{enumerate}

条件 a)、b)、c) 的等价性由命题 4 立即得出。当 E 是希尔伯特空间时，条件 a) 与 d) 的等价性由 V, p.~16 的推论 1 得出。

定义 4. --- 分离的准希尔伯特空间 E 中的标准正交且完全的族称为 E 的标准正交基。

分离的准希尔伯特空间 E 的标准正交基也是 E 的完备化的标准正交基。

设 \((e_{i})_{i\in I}\) 为 E 的标准正交基；对每个 \(x\in E\)，数 \(\langle e_{i}|x\rangle\) 滥用术语称为 x 关于基 \((e_{i})\) 的坐标。对 E 中每个 x 与 y，我们有

\[
\langle x | y \rangle = \sum_ {i \in \mathbf {I}} \overline {{\langle e _ {i} | x \rangle}} \langle e _ {i} | y \rangle .\tag{3}
\]

E 的标准正交基一般并不是 A, II, p.~25 中所定义的意义下 E 在域 K 上的基；为避免任何混淆，我们把上述引文意义下准希尔伯特空间 E 的基称为 E 在 K 上的代数基。

设 E 与 F 为两个分离的准希尔伯特空间，u 为从 E 到 F 中的连续线性映射。设 \((e_{i})_{i\in I}\)（相应地，\((f_{j})_{j\in J}\)）为 E（相应地，F）的标准正交基。令

\[
u _ {j i} = \langle f _ {j} | u (e _ {i}) \rangle
\]

其中 \(i \in \mathbf{I}\)，\(j \in \mathbf{J}\)。族 \((u_{ji})_{(i,j) \in \mathbf{I} \times \mathbf{J}}\) 称为 \(u\) 关于标准正交基 \((e_i)\) 与 \((f_j)\) 的矩阵。设 \(x \in \mathbf{E}\) 且 \(y = u(x)\)；若分别用 \(\xi_i = \langle e_i | x \rangle\) 与 \(\eta_j = \langle f_j | y \rangle\) 表示 \(x\) 与 \(y\) 的坐标，则 \(\eta_j = \sum_{i \in \mathbf{I}} u_{ji} \xi_i\) 对所有 \(j \in \mathbf{J}\) 成立。当 \((e_i)\) 是 \(E\) 的代数基且 \((f_j)\) 是 \(F\) 的代数基时，我们的定义与 A, II, § 10, No.~4 中的定义一致。

例. --- 设 E 为 R 上复值连续函数的空间，这些函数满足 \(f(x + n) = f(x)\)（\(x \in R\)，\(n \in Z\)）。我们赋予 E 由下式定义的内积

\[
\langle f | g \rangle = \int_ {0} ^ {1} \overline {{{{f (t)}}}} g (t) d t.
\]

于是 E 是分离的准希尔伯特空间，但不完备。对每个整数 \(n \in Z\)，令 \(e_{n}(x) = \mathbf{e}(nx)\)。立即可以看出族 \((e_{n})_{n \in \mathbf{Z}}\) 在 E 中标准正交。此外，E 上的一致收敛拓扑细于由范数 \(\|f\|_{2} = \langle f | f \rangle^{1/2}\) 得到的拓扑。族 \((e_{n})_{n \in \mathbf{Z}}\) 在 E 中对一致收敛是完全的（GT, X, § 4, No.~4），从而在准希尔伯特空间 E 中更是完全的。因此 \((e_{n})_{n \in \mathbf{Z}}\) 是 E 的标准正交基。

